\section{About the Committee}

\subsection{History and Establishment of UNESCO}

UNESCO, or the United Nations Educational, Scientific, and Cultural Organization, is a specialized UN agency which aims to strengthen international cooperation in areas such as education, science, culture, and information. UNESCO was created on November 16, 1945, after the devastation of World War II, with 37 nations signing its Constitution.\footnote{Meskell, L. (2018). \emph{A future in ruins: UNESCO, world heritage, and the dream of peace}. Oxford University Press.} It was founded on the belief that lasting peace must be built through intellectual and moral solidarity among people rather than relying solely on political agreements. This principle is reflected in the preamble of its founding charter, which states that “since wars begin in the minds of men, it is in the minds of men that the defences of peace must be constructed”.\footnote{United Nations. (1945). \emph{Charter of the United Nations}.\url{https://www.un.org/en/about-us/un-charter}} The organization works alongside the 194 Member States to address global challenges through initiatives that advance quality education, ensure access to reliable information, promote scientific cooperation and ethical standards for AI, and safeguard cultural heritage.\footnote{UNESCO. (2023). \emph{World Heritage List}.\url{https://whc.unesco.org/en/list/}}

\subsection{Mandate, Functions, and Role within the UN System}

As UNESCO is a specialised UN agency, it does not legislate like a State would and cannot force governments to implement specific policies.\footnote{Forrest, C. (2010). \emph{International law and the protection of cultural heritage}. Routledge.} The organisations’ recommendations,  operational guidelines, reports, declarations, and conventions are not legally binding, except in States where they have been ratified. However, UNESCO’s work does wield influence on Member States in the form of soft power, which is characterised by moral authority, global visibility, peer pressure, and funding coordination instead of a strict legal mandate, all while setting international norms and standards.\footnote{Meskell, L. (2018). \emph{A future in ruins: UNESCO, world heritage, and the dream of peace}. Oxford University Press.} Although these instruments depend on State approval and political will, they substantially influence national legislation, inform judicial interpretation, and provide reference points for public policy.\footnote{Forrest, C. (2010). \emph{International law and the protection of cultural heritage}. Routledge.} This provides delegates with freedom to newly define what constitutes a threat to cultural heritage especially in this new reality of global armed conflicts, where international actors are essential to the protection of both tangible and intangible heritage around the world.\footnote{Blake, J. (2000). On defining the cultural heritage. \emph{International and Comparative Law Quarterly, 49}(1), 61–85.\url{https://doi.org/10.1017/S002058930006396X}}

UNESCO also functions as a global observatory. Through bodies such as the UNESCO Institute for Statistics (UIS), the organization can collect cross-nationally comparable data, track the States’ progress in various fields, and consequently issue reports that help to better grasp what States lack and what needs to be changed on a larger scale. Sustainable Development Goal (SDG) 11.4 “Monitoring progress toward safeguarding the world’s cultural and natural heritage” collects data on per capita expenditure for heritage preservation across all levels of government and types of heritage.\footnote{United Nations. (2015). \emph{The 17 Sustainable Development Goals}.\url{https://sdgs.un.org/goals}} Through its Survey of Expenditure on Cultural and Natural Heritage, the UIS cooperates with States on the collection of reliable data to monitor global progress, identify areas of funding shortfalls, and ultimately promote the role of culture in sustainable development.\footnote{United Nations. (2025). A realistic goal? Culture in the Sustainable Development Agenda. \emph{UN Chronicle}.\url{https://www.un.org/en/un-chronicle/realistic-goal-culture-sustainable-development-agenda}} The UNESCO Framework for Cultural Statistics serves as the formalised system to measure and assess the social and economic impact of culture and the creative industries. The Framework provides governments, statistical institutions, and researchers with harmonised definitions, concepts, and methodologies that allow for developing internationally comparable data.

UNESCO is the worldwide leading coordinator and advocate for heritage protection. It brings together governments, international actors,  financial institutions, non-governmental organisations (NGOs), research institutions, and private actors to align their efforts around shared goals. As SDG 11 “Make cities and human settlements inclusive, safe, resilient, and sustainable”.

\subsection{UNESCO and the Protection of Cultural Heritage}

UNESCO considers cultural heritage to be more than a collection of objects and traditions.\footnote{Blake, J. (2000). On defining the cultural heritage. \emph{International and Comparative Law Quarterly, 49}(1), 61–85.\url{https://doi.org/10.1017/S002058930006396X}} Throughout its history and broad mandate, one of its most significant priorities has been safeguarding cultural heritage.\footnote{Meskell, L. (2018). \emph{A future in ruins: UNESCO, world heritage, and the dream of peace}. Oxford University Press.} The organization maintains that protecting the world’s cultural and natural heritage is key to addressing contemporary global challenges. This includes the increasing number of complex emergencies and armed conflicts.\footnote{UNESCO. (1972). \emph{Convention Concerning the Protection of the World Cultural and Natural Heritage}.\url{https://whc.unesco.org/en/convention/}} UNESCO has consistently emphasized that culture is not merely a consequence of growth but an essential driver to sustainable development, defined as a resolution that meets the needs of the present without compromising the future.

Cultural heritage has been central to UNESCO’s identity. Over the decades, the organization has shaped the core international treaties on heritage preservation and diversity. Such as the 1954 Hague Convention\footnote{UNESCO. (1954). \emph{Convention for the Protection of Cultural Property in the Event of Armed Conflict}.\url{https://en.unesco.org/protecting-heritage/convention-and-protocols/1954-convention}}, the 1972 World Heritage Convention\footnote{UNESCO. (1972). \emph{Convention Concerning the Protection of the World Cultural and Natural Heritage}.\url{https://whc.unesco.org/en/convention/}}, and the 2003 Convention for the Safeguarding of the Intangible Cultural Heritage,\footnote{UNESCO. (2003). \emph{Convention for the Safeguarding of the Intangible Cultural Heritage}.\url{https://ich.unesco.org/en/convention}} and many more. Together, these conventions provide UNESCO with a global legal framework to support Member States in the protection of cultural heritage. They encourage cooperation and dialogue between nations, create global standards for conservation, and provide the means to understand and protect culturally significant heritage.\footnote{Boylan, P. J. (1993). \emph{Review of the Convention for the Protection of Cultural Property in the Event of Armed Conflict (The Hague Convention of 1954)}. UNESCO.\url{https://unesdoc.unesco.org/ark:/48223/pf0000102816}}

In addition to the creation of international legal instruments, UNESCO’s commitment to protecting the cultural heritage of the world further extends through providing a larger contribution to the UN Sustainable Development Agenda. Today, it is the leading reference within the United Nations for the protection of both tangible and intangible cultural heritage and is acclaimed for its achievements in integrating culture as an essential part of sustainable development. This commitment is reflected in the United Nations 2030 Agenda through SDG 11.

\subsection{Funding}

UNESCO’s initiatives are constrained by a relatively modest and contribution-dependent budget. Its regular budget is composed primarily of assessed contributions from Member States, which covers core functions of the organization. Such as supporting the personnel and offices, both in the headquarters and in the field, as well as a minimum level of programme activities regarding heritage, education, science, culture, communication, etc. UNESCO’s abilities are more focused on contributing through its expertise, knowledge-sharing, coordination, and targeted capacity support rather than large-scale funding.

UNESCO relies heavily on voluntary contributions and extra-budgetary resources from States, other international organisations, foundations, and private partners to support many of its concrete initiatives.\footnote{Di Giovine, M. (2015). \emph{UNESCO's World Heritage Program: The challenges and ethics of community participation}. OpenEdition.\url{https://books.openedition.org/gup/213}} In this case, the organization relies on targeted trust funds and assistance programs, with a focus on emergency conservation, disaster recovery, and preparatory nomination assistance. The organisation provides annual World Heritage Fund activities, Heritage Emergency Fund (HEF), and open calls for proposals. As opposed to being strictly focused on providing funding, the organization is recognized as a resource for Member States when it comes to developing heritage protection strategies and securing financial assistance from national governments, development banks, multilateral organizations, and additional partners throughout the globe.

\subsection{Limitations and Challenges}

UNESCO has a broad mandate and considerable reach internationally. However, it faces numerous institutional and operational challenges. The organisation cannot enforce its decisions and must rely primarily on the independent implementation of Member States. UNESCO’s ability to respond to threats to heritage may also be limited by a lack of funds, inability to access conflict-impacted areas, and the complexity of coordinating international responses from numerous stakeholders.\footnote{European Parliament. (2023). \emph{Protecting cultural heritage from armed conflicts in Ukraine and beyond}. Policy Department for Structural and Cohesion Policies.\url{https://www.europarl.europa.eu/RegData/etudes/STUD/2023/733120/IPOL_STU\%282023\%29733120_EN.pdf}} Despite having an important role in developing international norms, the effective protection of cultural heritage is the obligation and responsibility of individual States and the international community.
